\section{A simultaneous pin test}\label{sec:pin-test}

We relate the probability that many specified points lie on cycles missing
the pins to a mean return probability after those points are pinned.
The estimate uses the complete pin law, before any trajectories are exposed.

Fix a finite volume and a finite period $\beta>0$.  Let $E$ be a finite
pin set, or the union of a finite pin set with $D_0$, and let $T$ be a
nonempty finite set of slot points at one time, disjoint from $E$.
Put $m=|T|$ and $P=E\cup T$.
For $i\in T$, let $\mathcal R_i(E,T)$ be the event that the line started
at $i$ reaches $T$ again, at strictly positive elapsed time, before it
reaches $E$.  The same sampled exchange picture is repeated when following
this line through successive periods.  Define
\begin{equation}\label{eq:pin-return-fraction}
 r=r(E,T):=\frac1m\sum_{i\in T}
                  \mathbb P_P\bigl(\mathcal R_i(E,T)\bigr).
\end{equation}
Thus $0\le r\le1$.  The positive elapsed-time convention excludes the
starting point from the return test.

\begin{proposition}[Simultaneous pin test]\label{prop:pin-test}
For every finite volume, every $\beta>0$, every pin set $E$ that is
finite or a finite set together with $D_0$, and every nonempty finite
single-time set $T$ disjoint from $E$, the return fraction in
\eqref{eq:pin-return-fraction} satisfies
\begin{equation}\label{eq:simultaneous-pin-test}
       \mathbb P_E(T\subseteq\mathcal B_E)\le (2r)^{|T|}.
\end{equation}
\end{proposition}

We first assume that $E$, and hence $P=E\cup T$, is finite.
For each exchange picture, let $\sigma$ be the permutation of $P$ that
maps each pin to the next pin on its directed cycle, at strictly positive
elapsed time.  The random subset $\sigma(T)$ has exactly $m$ elements.
Its relation to the event we want to bound is
\begin{equation}\label{eq:pin-event-equivalence}
       \{T\subseteq\mathcal B_E\}
             =\{\sigma(T)=T\}=\{T\subseteq\sigma(T)\}.
\end{equation}
Indeed, if every cycle meeting $T$ avoids $E$, its successive pins all
belong to $T$.  Conversely, if $\sigma(T)=T$, iterating $\sigma$ from
a point of $T$ visits every pin on that cycle without reaching $E$.
The last equality uses $|\sigma(T)|=|T|$.

The return fraction also has a subset interpretation:
\begin{equation}\label{eq:pin-return-output}
 mr=\mathbb E_P|T\cap\sigma(T)|
                  =\sum_{j\in T}\mathbb P_P(j\in\sigma(T)).
\end{equation}
Bijectivity of $\sigma$ makes the returning starting points and their
outputs in $T$ have the same cardinality.  We will prove stability
of the subset generating polynomial
$\mathbb E_P\prod_{j\in\sigma(T)}y_j$ and deduce that the probability
of containing all of $T$ is at most the product of its presence
probabilities.  By \eqref{eq:pin-return-output} that product is at most
$r^m$.  A final comparison between the laws $\mathbb P_E$ and
$\mathbb P_P$ supplies the factor $2^m$ in
\eqref{eq:simultaneous-pin-test}.

\subsection{The stability operations}

Write $\mathbb H=\{z\in\C:\operatorname{Im}z>0\}$.
A nonzero polynomial in finitely many complex variables is
\emph{stable} if it has no zero when every variable belongs to
$\mathbb H$.  A polynomial is \emph{multiaffine} if its degree in
each variable separately is at most one.

\begin{lemma}[Closure operations]\label{lem:pin-stability-closure}
For polynomials in a fixed finite set of variables the following
operations preserve stability, with the possibility of producing the
zero polynomial: coefficientwise limits with a common degree bound,
specialization of a variable to a real number, and partial differentiation.
Identifying two variables preserves stability.  In particular, identifying
two variables as $t$ and extracting the coefficient of $t$ preserves
stability up to zero.
\end{lemma}

\begin{proof}
Let $F_n$ be stable polynomials of bounded degree converging coefficientwise
to a nonzero polynomial $F$.  If $F$ vanished at a point of a product of
upper half-planes, choose a complex line through that point on which the
restriction of $F$ is not identically zero.  A sufficiently small disk
in this line lies in the product of upper half-planes.  The restrictions
of $F_n$ converge uniformly on the disk.  On a smaller circle avoiding
the isolated zeros of the restriction of $F$, the argument principle
(or Rouch\'e's theorem) forces $F_n$ to have a zero inside for large $n$.
This is impossible.  The existence of the line follows, for example, by
choosing a direction where the first nonzero homogeneous term of the
Taylor expansion of $F$ at the purported zero does not vanish.

Specialization at $a\in\R$ is the limit of specialization at
$a+\mathrm i/n$, so the first assertion applies.  Identification of
variables is immediate from the definition.

For differentiation in $t$, let $k$ be the degree of $F$ in $t$.
The case $k=0$ gives the zero polynomial.  If $k\ge1$, the leading
coefficient is the nonzero coefficientwise limit
\[
 \lim_{R\to\infty}(\mathrm iR)^{-k}
                       F(z_1,\ldots,z_{n-1},\mathrm iR).
\]
It is therefore stable.  Consequently, after fixing all other variables
in $\mathbb H$, the resulting polynomial in $t$ still has degree $k$.
Write its roots, with multiplicity, as $\zeta_1,\ldots,\zeta_k$;
each has imaginary part at most zero.  For $t\in\mathbb H$,
\[
 \frac{\partial_t F}{F}
       =\sum_{j=1}^k\frac1{t-\zeta_j}
\]
has strictly negative imaginary part.  It cannot vanish, proving
stability of the derivative.  Finally, taking the linear coefficient
after identification is differentiation in $t$ followed by $t=0$.
\end{proof}

The next operation is the partial-symmetrization principle in the
stable-polynomial approach to negative dependence
\cite[Theorem~4.20]{BorceaBrandenLiggett2009}.  For an elementary proof
using complex bivariate sections, see also
Liggett~\cite[Section~5]{Liggett2009Distributional}.  We give the
calculation in the form needed here.

\begin{lemma}[Partial symmetrization]\label{lem:pin-partial-symmetrization}
Let $F$ be a stable multiaffine polynomial, let $\tau$ exchange two
of its variables, and let $0\le u\le1$.  Then
$(1-u)F+u\tau F$ is stable and multiaffine.
\end{lemma}

\begin{proof}
Fix every other variable in $\mathbb H$ and write the two-variable
section as
\[
                 q(z,w)=a+bz+cw+ezw,
\]
where $a,b,c,e$ can be complex.  For every $z\in\mathbb H$,
\begin{equation}\label{eq:pin-imaginary-criterion}
 N(z):=\operatorname{Im}\bigl((a+bz)\overline{(c+ez)}\bigr)\ge0.
\end{equation}
Indeed, if $c+ez\ne0$, the root $-(a+bz)/(c+ez)$ in the $w$
variable cannot belong to $\mathbb H$; if $c+ez=0$, the displayed
quantity is zero.  The swapped section $q(w,z)$ is also stable, so
\[
 N^\tau(z):=\operatorname{Im}\bigl((a+cz)\overline{(b+ez)}\bigr)\ge0.
\]
Put $b_u=(1-u)b+uc$ and $c_u=(1-u)c+ub$.  Direct expansion gives
\begin{align}\label{eq:pin-symmetrization-identity}
 \operatorname{Im}\bigl((a+b_uz)\overline{(c_u+ez)}\bigr)
  &=(1-u)N(z)+uN^\tau(z)
       +u(1-u)|b-c|^2\operatorname{Im}z.
\end{align}
If $0<u<1$ and $b\ne c$, the right-hand side is strictly positive.
In particular $c_u+ez\ne0$, and the root in $w$ of the mixed section
lies strictly below the real axis.  If $u\in\{0,1\}$ or $b=c$,
the section is the original section or its swap.  These cases also
exclude zeros in $\mathbb H^2$.  Since the other variables were arbitrary,
the mixed polynomial is stable; multiaffinity is preserved directly.
\end{proof}

The complex coefficients in this proof matter: other variables have
already been fixed in the upper half-plane.  No reality assumption on
those sections was used.

\subsection{Layering the exchange picture and removing unpinned points}

For a permutation $\pi$ of a finite set $V$, define its polynomial
\[
             F_\pi(\boldsymbol x,\boldsymbol y)
                     :=\prod_{i\in V}(x_i+y_{\pi(i)}).
\]
Each $x_i$ and each $y_i$ occurs with degree at most one.

\begin{lemma}[Stable exchange symbols]\label{lem:pin-exchange-symbol}
For a permutation obtained by a finite product of independent Bernoulli
transpositions and deterministic permutations on a finite set $V$,
the polynomial $\mathbb E F_\pi$ is stable.  Here each Bernoulli
transposition uses an arbitrary probability in $[0,1]$.  The same
conclusion holds for any limit of these permutation laws.

Every finite independent product of finite-time Poisson exchange
permutations, uniform permutations on specified subsets of $V$, and
deterministic permutations has such a limit representation.
\end{lemma}

\begin{proof}
For the identity permutation the polynomial is
$\prod_{i\in V}(x_i+y_i)$, a product of nonvanishing upper-half-plane
sums.  Composing the output with a transposition either leaves the
polynomial unchanged or interchanges the corresponding two $y$ variables.
Averaging an independent Bernoulli choice is precisely the operation in
Lemma~\ref{lem:pin-partial-symmetrization}.  Deterministic permutations
only relabel variables.  Repetition proves the first assertion.
Convergence of permutation laws gives coefficientwise convergence of
their symbols, whose degrees are bounded and whose value at all variables
equal to $1$ is $2^{|V|}$.  Lemma~\ref{lem:pin-stability-closure} therefore
proves stability of every such limit, with zero excluded.

For a finite-time Poisson exchange process, partition its time interval
into slices of length at most $\delta$.  On each edge and slice record
whether the Poisson clock has at least one mark; these indicators are
independent Bernoulli variables.  Apply the indicated transpositions
in a fixed order within a slice.  This coupled approximation agrees with
the actual process whenever every slice has at most one mark in total.
If the total rate is $\lambda$ and the interval length is $t$, the
probability of disagreement is at most $C\lambda^2(t+\delta)\delta$,
which tends to zero.  This constructs the required limit of
Bernoulli-product laws.

To approximate a uniform permutation on a subset $A\subseteq V$,
run Poisson transpositions at positive rates on every pair in $A$.
For $|A|\le1$ there is nothing to prove.  Otherwise the unit-time
transition matrix on the finite group of permutations of $A$ has all
entries positive: every target permutation is a finite word in these
transpositions, and that word can occur in the prescribed order within
unit time with no other marks.  The uniform distribution is invariant,
since each step is multiplication by a random group element.  If the
group has $M$ elements and the minimum transition entry is $\eta>0$,
the transition matrix has a common uniform part of mass $M\eta$.
Splitting off that part shows that distance from the uniform law
contracts by at most $1-M\eta$ at each unit step (and is already zero
if $M\eta=1$).  Hence the law converges to uniform as time tends to
infinity.  At each finite time the small-slice construction gives a
Bernoulli-product approximation.  Taking a sequence with both errors
tending to zero proves the asserted representation of the uniform law.
For a finite independent product, approximate its factors independently;
their product laws then converge to the desired product law.
\end{proof}

Now assume $P=E\cup T$ is finite.  Cut the time circle at every time
represented in $P$, and include the time $0$ cut if necessary.  Let
$V$ be the disjoint union of a copy of $V_\Lambda$ at every cut.
Assign an endpoint permutation to the interval ending at time $0$,
so that this cut has the post-completion convention.  Following a line
from one cut to the next defines a permutation $\pi$ of $V$.  Its
cycles are exactly the cycles of the exchange picture, with the cut
points inserted as vertices.

The random interval bijections are independent.  Moreover, the $x$
variables on an interval's starting layer and the $y$ variables on its
ending layer are disjoint from the variable sets used by all other
intervals.  Consequently
\begin{equation}\label{eq:pin-layered-symbol}
 F_V(\boldsymbol x,\boldsymbol y)
       :=\mathbb E_{\mathrm b}\prod_{i\in V}(x_i+y_{\pi(i)})
\end{equation}
is the product of the stable interval symbols of
Lemma~\ref{lem:pin-exchange-symbol}, with cyclic relabeling of their
output variables.  It is therefore stable.  This also covers a single
cut, where there is just one interval wrapping around the circle.

The next-pin permutation $\sigma$ is obtained from this layered
permutation by following its cycles between successive points of $P$.
Let $n_P(\pi)$ denote the number of cycles of $\pi$ missing $P$.

\begin{lemma}[Splicing and output extraction]\label{lem:pin-splicing}
Starting from \eqref{eq:pin-layered-symbol}, eliminate each $q\in V\setminus P$
by setting $x_q=y_q=t$ and extracting the coefficient of $t$.
The resulting nonzero stable polynomial is
\begin{equation}\label{eq:pin-spliced-symbol}
 F_P(\boldsymbol x,\boldsymbol y)
   =\mathbb E_{\mathrm b}
       \left[2^{n_P(\pi)}\prod_{i\in P}(x_i+y_{\sigma(i)})\right].
\end{equation}
If $T\subseteq P$, then the probability generating polynomial of the
random subset $\sigma(T)$ under $\mathbb P_P$ is stable.  More precisely,
with $Z_P=\mathbb E_{\mathrm b}2^{n_P(\pi)}>0$, it is
\begin{equation}\label{eq:pin-all-x-zero}
 G_T(\boldsymbol y)
  :=\mathbb E_P\prod_{i\in T}y_{\sigma(i)}
   =\frac1{Z_P}
     \left.
       \left(\prod_{i\in P\setminus T}\partial_{x_i}\right)
       F_P(\boldsymbol x,\boldsymbol y)
     \right|_{x_i=0\ (i\in P)}.
\end{equation}
\end{lemma}

\begin{proof}
Consider first one deterministic permutation.  If $q$ is not a fixed
point, write $i=\pi^{-1}(q)$ and $j=\pi(q)$; the two factors involving
$y_q$ and $x_q$ become
\[
             (x_i+t)(t+y_j).
\]
Their linear coefficient is $x_i+y_j$.  Thus deleting $q$ splices
the directed edges $i\to q\to j$ into $i\to j$.  This includes a
two-cycle, for which $i=j$.  If $q$ is a fixed point, its factor is
$x_q+y_q=2t$, giving a factor $2$.  Successive deletion leaves the
next-pin permutation on every cycle meeting $P$.  A cycle missing
$P$ eventually becomes a fixed point and contributes exactly one
factor $2$.  Linearity gives \eqref{eq:pin-spliced-symbol};
Lemma~\ref{lem:pin-stability-closure} gives stability up to zero.
Its nonnegative coefficients and positive value at all surviving
variables equal to $1$ exclude zero.

In each product on the right of \eqref{eq:pin-spliced-symbol},
differentiation in $x_i$, $i\in P\setminus T$, selects the $x_i$
term from exactly those factors.  Setting \emph{every} remaining $x$
variable to zero, including those indexed by $T$, selects $y_{\sigma(i)}$
from each remaining factor.  Division by $Z_P$ gives
\eqref{eq:pin-all-x-zero}.  Derivatives and real specialization preserve
stability up to zero; $G_T(\boldsymbol1)=1$ excludes zero.  Since
$\sigma$ is a permutation, the product contains $m=|T|$ distinct
$y$ variables, so this is indeed a multiaffine subset generating
polynomial with nonnegative coefficients.
\end{proof}

We have now encoded the outputs of $T$ under the correct pin weight.
The next lemma turns stability of this particular generating polynomial
into the simultaneous return bound.  This is the subset-presence
consequence of the strongly Rayleigh framework
\cite[Theorem~4.9]{BorceaBrandenLiggett2009}; the short conditional
covariance argument below proves exactly what is required.

\subsection{Covariance, normalization, and the dense-boundary limit}

\begin{lemma}[Presence probabilities]\label{lem:pin-presence}
Let $A$ be a random subset of a finite set $I$.  Suppose its generating
polynomial $G(\boldsymbol y)=\mathbb E\prod_{j\in A}y_j$ is stable.
For every $J\subseteq I$,
\begin{equation}\label{eq:pin-presence-product}
           \Pr(J\subseteq A)\le\prod_{j\in J}\Pr(j\in A).
\end{equation}
The same assertion holds after conditioning on $B\subseteq A$, whenever
that event has positive probability, for subsets of $I\setminus B$.
\end{lemma}

\begin{proof}
Write $X_j=\1_{\{j\in A\}}$.  Conditioning on $B\subseteq A$ gives,
for the indicators outside $B$, the generating polynomial
\[
 \frac{\left.\left(\prod_{j\in B}\partial_{y_j}\right)
                       G(\boldsymbol y)\right|_{y_j=1\ (j\in B)}}
      {\Pr(B\subseteq A)}.
\]
It is nonzero and stable by Lemma~\ref{lem:pin-stability-closure}.
Under this conditional law, keep only two variables $z,w$ and
specialize all others to $1$.  The resulting stable polynomial has
the form $a+bz+cw+ezw$, where $a,b,c,e\ge0$ and
$a+b+c+e=1$.  Applying \eqref{eq:pin-imaginary-criterion} with real
coefficients gives
\[
   0\le\operatorname{Im}\bigl((a+bz)\overline{(c+ez)}\bigr)
                 =(bc-ae)\operatorname{Im}z.
\]
Thus $bc\ge ae$.  The covariance of the corresponding presence
indicators equals
\[
                  e-(b+e)(c+e)=ae-bc\le0.
\]
Whenever the second indicator has positive probability of being one,
conditioning it to be one therefore cannot increase the probability
that the first is one.

Order $J=\{j_1,\ldots,j_k\}$.  If any prefix-presence event has
probability zero, \eqref{eq:pin-presence-product} is immediate.
Otherwise, for each $h$, remove the conditions
$X_{j_{h-1}}=1,\ldots,X_{j_1}=1$ one at a time.  The preceding
conditional covariance inequality gives
\[
 \Pr\bigl(X_{j_h}=1\mid X_{j_1}=\cdots=X_{j_{h-1}}=1\bigr)
                        \le \Pr(X_{j_h}=1).
\]
Multiplying the conditional probabilities proves
\eqref{eq:pin-presence-product}.  Applying the same argument to the
conditional generating polynomial proves the final assertion.
\end{proof}

\begin{proof}[Proof of Proposition~\ref{prop:pin-test}]
Suppose first that $E$ is finite, and use the layered construction
above with $P=E\cup T$.  By
Lemmas~\ref{lem:pin-splicing} and~\ref{lem:pin-presence}, followed
by the arithmetic--geometric mean inequality and
\eqref{eq:pin-return-output},
\begin{equation}\label{eq:pin-return-product}
 \mathbb P_P\bigl(T\subseteq\sigma(T)\bigr)
       \le\prod_{j\in T}\mathbb P_P(j\in\sigma(T))\le r^m.
\end{equation}
If one of the factors is zero this conclusion is interpreted directly;
no conditioning on a zero-probability event is needed.

By \eqref{eq:pin-event-equivalence}, the event on the left is
$H=\{T\subseteq\mathcal B_E\}$.  It remains to compare its probability
under the two pin laws.

For clarity, write $w_D=2^{n_D}$, where $n_D$ is the number of cycles
missing $D$, and $Z_D=\mathbb E_{\mathrm b}w_D$.
Deleting the $m$ pins in $T$ changes the number of unpinned cycles by
the number of cycles that meet $T$ but miss $E$.  That number lies
between $0$ and $m$.  Hence, picture by picture,
\[
                w_P\le w_E\le 2^m w_P,
                \qquad Z_E\ge Z_P>0.
\]
We obtain
\[
 \mathbb P_E(H)
  =\frac{\mathbb E_{\mathrm b}(\1_Hw_E)}{Z_E}
  \le 2^m\frac{Z_P}{Z_E}\mathbb P_P(H)
  \le 2^m r^m.
\]
This proves the finite-pin case with the required normalization.

Finally let $E=F\cup D_0$ with $F$ finite.  Choose increasing finite
sets of boundary cuts with dense union, and let $D_{0,n}$ contain all
boundary slots at those cuts.  Put $E_n=F\cup D_{0,n}$ and
$P_n=E_n\cup T$.  Then $E_n$ and $T$ are disjoint and the finite-pin
inequality applies.  As in Lemma~\ref{lem:dense-pins}, almost every
picture has finitely many marks, and every boundary visit contains
an interval of positive duration.  Thus every boundary-touching
cycle is eventually detected, and
\[
       w_{E_n}\longrightarrow w_E,\qquad
       w_{P_n}\longrightarrow w_P,\qquad
       \1_{\{T\subseteq\mathcal B_{E_n}\}}
              \longrightarrow\1_{\{T\subseteq\mathcal B_E\}}.
\]
The weights are bounded by $2^{|V_\Lambda|}$ in this fixed volume,
and all their normalizers are positive, so the corresponding
normalized probabilities converge by dominated convergence.

The same reasoning applies to the return test, with a useful finite
horizon.  Starting from $i\in T$, follow the uncensored line until
its first positive return to $T$.  This time is finite, since the
line eventually returns to $i$ on its finite permutation cycle.
If the line meets $F$ before that return, every finite-cut test fails.
If it first meets the boundary before that return, one of its
positive-duration boundary visits before the return contains a cut
from the dense union, so all sufficiently fine tests fail.  If it
meets neither, every test succeeds.  Since $T$ is disjoint from
$D_0$, a boundary visit cannot first occur only at the return point.
Consequently, for every $i\in T$,
\[
 \1_{\mathcal R_i(E_n,T)}\longrightarrow
             \1_{\mathcal R_i(E,T)}
             \quad\text{almost surely}.
\]
Bounded convergence with the weights $w_{P_n}$ now gives
$r(E_n,T)\to r(E,T)$.  Taking limits in
$\mathbb P_{E_n}(T\subseteq\mathcal B_{E_n})\le(2r(E_n,T))^m$
proves \eqref{eq:simultaneous-pin-test} for continuous boundary pins.
Only scalar probabilities and bounded weights are passed to this
limit; no polynomial with infinitely many variables is required.
\end{proof}
